% Standalone problem document, generated from the adjacent README.md.
% Compile with XeLaTeX. Mathematical content is maintained in Markdown.
\documentclass[11pt,a4paper]{article}
\usepackage[fontsize=10.5pt]{scrextend}
\usepackage[margin=22mm,headheight=15pt,headsep=9mm,footskip=11mm]{geometry}
\usepackage{fontspec}
\setmainfont{texgyrepagella}[Extension=.otf,UprightFont=*-regular,BoldFont=*-bold,ItalicFont=*-italic,BoldItalicFont=*-bolditalic]
\setsansfont{texgyreheros}[Extension=.otf,UprightFont=*-regular,BoldFont=*-bold,ItalicFont=*-italic,BoldItalicFont=*-bolditalic]
\setmonofont{texgyrecursor}[Extension=.otf,UprightFont=*-regular,BoldFont=*-bold,ItalicFont=*-italic,BoldItalicFont=*-bolditalic,Scale=MatchLowercase]
\usepackage{amsmath,amssymb}
\usepackage{unicode-math}
\setmathfont{texgyrepagella-math.otf}
\usepackage{xcolor,microtype,enumitem,needspace,fancyhdr}
\usepackage{bookmark}
\definecolor{ink}{HTML}{17344B}
\definecolor{accent}{HTML}{197D80}
\definecolor{muted}{HTML}{596773}
\hypersetup{colorlinks=true,linkcolor=accent,urlcolor=accent,pdftitle={IS-05 - The
optimal decay exponent for the conditioning of sign
matrices},pdfauthor={Open Problems in Numerical Linear Algebra}}
\urlstyle{same}
\setlength{\parindent}{0pt}
\setlength{\parskip}{5pt plus 1pt minus 1pt}
\linespread{1.04}
\setlength{\emergencystretch}{3em}
\widowpenalty=10000
\clubpenalty=10000
\displaywidowpenalty=10000
\allowdisplaybreaks[1]
\setlist{nosep,leftmargin=1.5em}
\providecommand{\tightlist}{\setlength{\itemsep}{0pt}\setlength{\parskip}{0pt}}
\providecommand{\pandocbounded}[1]{#1}
\pagestyle{fancy}
\fancyhf{}
\fancyhead[L]{\sffamily\footnotesize\color{muted}OPEN PROBLEMS IN NUMERICAL LINEAR ALGEBRA}
\fancyhead[R]{\sffamily\bfseries\footnotesize\color{accent}IS-05}
\fancyfoot[L]{\parbox[b]{0.9\textwidth}{\sffamily\fontsize{8}{10}\selectfont\color{muted}Editorial ratings; a bounded search cannot certify that a problem remains unsolved.\\Literature check: 2026-09-08}}
\fancyfoot[R]{\sffamily\footnotesize\color{muted}\thepage}
\renewcommand{\headrulewidth}{0.3pt}
\renewcommand{\headrule}{\hbox to\headwidth{\color{accent}\leaders\hrule height \headrulewidth\hfill}}
\makeatletter
\renewcommand{\section}{\Needspace{5\baselineskip}\@startsection{section}{1}{0pt}{12pt plus 2pt minus 2pt}{4pt}{\sffamily\bfseries\large\color{ink}}}
\renewcommand{\subsection}{\Needspace{4\baselineskip}\@startsection{subsection}{2}{0pt}{9pt plus 1pt minus 1pt}{3pt}{\sffamily\bfseries\normalsize\color{ink}}}
\makeatother
\setcounter{secnumdepth}{0}
\begin{document}
{\sffamily\small\color{muted}Eigenvalues and inverse problems\par}
\vspace{3pt}
{\raggedright\hyphenpenalty=10000\exhyphenpenalty=10000\sffamily\bfseries\fontsize{21}{24}\selectfont\color{ink}The
optimal decay exponent for the conditioning of sign matrices\par}
\vspace{7pt}
{\sffamily\small\textbf{Difficulty:} extreme\quad\textbf{Importance:} interesting
to the community\par}
\vspace{4pt}
{\color{accent}\hrule height 0.8pt}
\vspace{6pt}
\textbf{Status:} open.

\subsection{Context and notation}\label{context-and-notation}

For \(n\geq1\), define

\[
h(n)=\min_{A\in\{-1,1\}^{n\times n}}\kappa_2(A),\qquad
\kappa_2(A)=\frac{\sigma_{\max}(A)}{\sigma_{\min}(A)},
\]

with \(\kappa_2(A)=+\infty\) for singular \(A\).

\subsection{Problem statement}\label{problem-statement}

Determine the exact number

\[
\alpha_*=\sup\{\alpha\geq0:\ \exists C>0\ \forall n\geq1,
\ h(n)-1\leq Cn^{-\alpha}\}.
\]

The constant \(C\) may depend on \(\alpha\), but not on dimension. This
specifies the uniform asymptotic power exponent; logarithmic factors do
not change the supremum. Orders admitting exact Hadamard matrices have
\(h(n)-1=0\), which are included without taking logarithms of zero.
Current bounds give \(17/92\leq\alpha_*\leq1\).

\subsection{References}\label{references}

Alexeev, Jasper, and Mixon,
\href{https://arxiv.org/html/2511.14653v1}{\emph{Asymptotically optimal
approximate Hadamard matrices}}, §6, Problem 11; the supremum above is
an editorial precise formulation of its decay-exponent question.
Steinerberger,
\href{https://faculty.washington.edu/steinerb/openproblems.pdf}{\emph{Open
Problems}}, Problem 69, November 2025 update, identifies the sharp rate
as unresolved.

\subsection{Status check}\label{status-check}

Searches for \textit{approximate Hadamard 17/92 2026} and
\textit{approximate Hadamard condition 2026 sharp} found no exact
exponent. A global constant as in
\href{https://github.com/ajt60gaibb/OpenProblemsInNLA/blob/main/eigenvalues-and-inverse-problems/IS-04/README.md}{IS-04}
does not specify a decay exponent, while an asymptotic exponent permits
finitely many exceptions to any proposed sharp constant. The two
independently posed problems are therefore retained separately.
\end{document}
